\Chapter{43}

\Verse{1}
{
    \zA{话说王夫人因见贾母那日在大观园不过着了些风寒，不是什么大病，请医生吃 了两剂药也就好了，命凤姐来，吩咐他预备给贾政带送东西。 正商议着，只见贾母
        打发人来叫，王夫人忙引着凤姐儿过来。 王夫人又请问：“这会子可又觉大安些？” 贾母道：“今日可大好了。 方才你们送来野鸡崽子汤，我尝了一尝，倒有味儿，又
        吃了两块肉，心里很受用。” 王夫人笑道：“这是凤丫头孝敬老太太的，算他的孝心 虔，不枉了素日老太太疼他。” 贾母点头笑道：“难为他想着。}
}
{
    \eA{When Madame Wang saw, for we will now proceed with our narrative, that the
        extent of dowager lady Chia's indisposition, contracted on the day she had
        been into the garden of Broad Vista, amounted to a simple chill, that no
        serious ailment had supervened, and that her health had improved soon after
        the doctor had been sent for and she had taken a couple of doses of
        medicine, she called lady Feng to her and asked her to get ready a present
        of some kind for her to take to her husband, Chia Cheng.}
}
{
}

\Verse{2}
{
    \zA{若是还有生的，再炸 上两块，咸浸浸的，喝粥有味儿。 那汤虽好，就只不对稀饭。” 凤姐听了，连忙答 应，命人到大厨房传话。
        这里贾母又向王夫人笑道：“我打发人找你来，不为别的：初二日是凤丫头的 生日。 上两年我原想着替他做生日，偏到跟前又有事就混过去了。 今年人又齐全，
        料着又没事，咱们大家好生乐一天。” 王夫人笑道：“我也想着呢。 既是老太太高兴， 何不就商议定了？”}
}
{
    \eA{But while they were engaged in deliberation, they perceived a waiting-maid
        arrive. She came from their old senior's part to invite them to go to her.
        So, with speedy step, Madame Wang led the way for lady Feng, and they came
        over into her quarters. "Pray, may I ask," Madame Wang then inquired,
        "whether you're feeling nearly well again now?" "I'm quite all right
        to-day," old lady Chia replied.}
}
{
}

\Verse{3}
{
    \zA{贾母笑道：“我想往年不拘谁做生日，都是各自送各自的礼， 这个也俗了，也觉太生分。 今儿我出个新法子，又不生分，又可以取乐儿。” 王夫
        人忙道：“老太太怎么想着好，就是怎么样行。” 贾母笑道：“我想着咱们也学那小 家子，大家凑个分子，多少尽着这钱去办，你说好不好？”
        王夫人道：“这个很好， 但不知怎么个凑法儿？”}
}
{
    \eA{"I've tasted the young-pheasant soup you sent me a little time back and find
        it full of relish. I've also had two pieces of meat, so I feel quite
        comfortable within me." "These dainties were presented to you, dear
        ancestor, by that girl Feng," Madame Wang smiled. "It only shows how sincere
        her filial piety is. She does not render futile the love, which you,
        venerable senior, ever lavish on her."}
}
{
}

\Verse{4}
{
    \zA{贾母听说，一发高兴起来，忙遣人去请薛姨妈邢夫人等， 又叫请姑娘们并宝玉，和那府里的尤氏和赖大家的，及有些头脸管事的媳妇也都叫 了来。
        众丫头婆子见贾母十分高兴，也都高兴，忙忙的各自分头去请的请，传的传。 没顿饭的工夫，老的少的，上的下的，乌压压挤了一屋子。 只薛姨妈和贾母对坐，
        邢夫人王夫人只坐在房门前两张椅子上，宝钗姐妹等五六个人坐在炕上，宝玉坐在 贾母怀前，底下满满的站了一地。}
}
{
    \eA{Dowager lady Chia nodded her head assentingly. "She's too kind to think of
        me!" she answered smiling. "But should there be any more uncooked, let them
        fry a couple of pieces; and, if these be thoroughly immersed in wine, the
        congee will taste well with them. The soup is, it's true, good, but it
        shouldn't, properly speaking, be prepared with fine rice."}
}
{
}

\Verse{5}
{
    \zA{贾母忙命拿几张小杌子来，给赖大母亲等几个高 年有体面的嬷嬷坐了。 贾府风俗：年高伏侍过父母的家人比年轻的主子还有体面呢，
        所以尤氏凤姐等只管地下站着，那赖大的母亲等三四个老嬷嬷告了罪，都坐在小杌 子上。 贾母笑着把方才一夕话说与众人听了，众人谁不凑这趣儿呢。 再也有和凤姐儿
        好，情愿这样的。 也有怕凤姐儿，巴不得奉承他的。 况且都是拿的出来的，所以一 闻此言都欣然应诺。 贾母先道：“我出二十两。”}
}
{
    \eA{After listening to her wishes, lady Feng expressed with alacrity her
        readiness to see them executed, and directed a servant to go and deliver the
        message in the cook-house. "I sent the servant for you," dowager lady Chia
        meanwhile said to Madame Wang with a smile, "not for anything else, but for
        the birthday of that girl Feng, which falls on the second.}
}
{
}

\Verse{6}
{
    \zA{薛姨妈笑道：“我随着老太太，也 是二十两。” 邢夫人王夫人笑道：“我们不敢和老太太并肩，自然矮一等，每人十六 两罢了。”
        尤氏李纨也笑道：“我们自然又矮一等，每人十二两罢。” 贾母忙和李纨 道：“你寡妇失业的，那里还拉你出这个钱，我替你出了罢。” 凤姐忙笑道：“老太
        太别高兴，且算一算账再揽事。 老太太身上已有两分呢。 这会子又替大嫂子出十二 两，说着高兴，一会子回想又心疼了!过后儿又说：‘都是为凤丫头花了钱。’}
}
{
    \eA{I had made up my mind two years ago to celebrate her birthday in proper
        style, but when the time came, there happened to be again something
        important to attend to, and it went by without anything being done. But this
        year, the inmates are, on one hand, all here, and there won't, I fancy, be,
        on the other, anything to prevent us, so we should all do our best to enjoy
        ourselves thoroughly for a day."}
}
{
}

\Verse{7}
{
    \zA{使个 巧法子，哄着我拿出三四倍子来暗里补上，我还做梦呢！” 说的众人都笑了。 贾母 笑道：“依你怎么样呢？”
        凤姐笑道：“生日没到，我这会子已经折受的不受用了。 我一个钱也不出，惊动这些人，实在不安，不如大嫂子这分我替他出了罢。 我到那
        一日多吃些东西，就享了福了。” 邢夫人等听了，都说很是，贾母方允了。 凤姐儿又笑道：“我还有一句话呢：我想老祖宗自己二十两，又有林妹妹宝兄 弟的两分子；
        姨妈自己二十两，又有宝妹妹的一分子：这倒也公道。}
}
{
    \eA{"I was thinking the same thing," Madame Wang rejoined, laughingly, "and,
        since it's your good pleasure, venerable senior, why, shouldn't we
        deliberate at once and decide upon something?" "To the best of my
        recollection," dowager lady Chia resumed smiling, "whenever in past years
        I've had any birthday celebrations for any one of us, no matter who it was,
        we have ever individually sent our respective presents;}
}
{
}

\Verse{8}
{
    \zA{只是二位太太 每位十六两，自己又少，又不替人出，这有些不公道。 老祖宗吃了亏了！” 贾母听 了，呵呵大笑道：“到底是我的凤丫头向着我，这说的很是。
        要不是你，我叫他们 又哄了去了。” 凤姐笑道：“老祖宗只把他哥儿两个交给两位太太，一位占一个罢， 派每位替出一分就是了。”
        贾母忙说：“这很公道，就是这样。”}
}
{
    \eA{but this method is common and is also apt, I think, to look very much as if
        there were some disunion. But I'll now devise a new way; a way, which won't
        have the effect of creating any discord, and will be productive of good
        cheer." "Let whatever way you may think best, dear ancestor, be adopted."
        Madame Wang eagerly rejoined.}
}
{
}

\Verse{9}
{
    \zA{赖大的母亲忙站起 来笑道：“这可反了，我替二位太太生气!在那边是儿子媳妇，在这边是内侄女儿，
        倒不向着婆婆姑姑，倒向着别人，这儿媳妇倒成了陌路人，‘内’侄女儿倒成了‘外’ 侄女儿了！” 说的贾母和众人都大笑起来了。
        赖大的母亲因又问道：“少奶奶们十二两，我们自然也该矮一等了？” 贾母听 说，道：“这使不得。 你们虽该矮一等，我知道你们这几个都是财主，位虽低些，
        钱却比他们多。 你们和他们一例才使得。” 众嬷嬷听了，连忙答应。}
}
{
    \eA{"My idea is," old lady Chia laughingly continued, "that we too should follow
        the example of those poor families and raise a subscription among ourselves,
        and devote the whole of whatever we may collect to meet the outlay for the
        necessary preparations. What do you say, will this do or not?" "This is a
        splendid idea!" Madame Wang acquiesced. "But what will, I wonder, be the way
        adopted for raising contributions?"}
}
{
}

\Verse{10}
{
    \zA{贾母又道：“姑 娘们不过应个景儿，每人照一个月的月例就是了。” 又回头叫鸳鸯来：“你们也凑几 个人，商议凑了来。”
        鸳鸯答应着，去不多时，带了平儿、袭人、彩霞等，还有几 个丫头来，也有二两的，也有一两的。 贾母因问平儿：“你难道不替你主子做生日？ 还入在这里头？”
        平儿笑道：“我那个私自另外的有了，这是公中的，也该出一分。” 贾母笑道：“这才是好孩子。” 凤姐又笑道：“上下都全了；
        还有二位姨奶奶，他出不出也问一声儿。}
}
{
    \eA{Old lady Chia was the more inspirited by her reply. There and then she
        despatched servants to go and invite Mrs. Hsüeh, Madame Hsing and the rest
        of the ladies, and bade others summon the young ladies and Pao-yü. But from
        the other mansion, Chia Chen's spouse, Lai Ta's wife, even up to the wives
        of such stewards as enjoyed a certain amount of respectability, were
        likewise to be asked to come round.}
}
{
}

\Verse{11}
{
    \zA{尽到他 们是理，不然他们只当小看了他们了。” 贾母听说，忙说：“可是呢。 怎么倒忘了他 们？ 只怕他们不得闲儿，叫个丫头问问去。” 说着，早有丫头去了。
        半日回来说道： “每位也出二两。” 贾母喜欢道：“拿笔砚来算明，共计多少。” 尤氏因悄悄的骂凤
        姐道：“我把你这没足够的小蹄子儿!这么些婆婆婶子凑银子给你做生日，你还不 够，又拉上两个苦瓠子。” 凤姐也悄悄的笑道：“你少胡说，一会子离了这里，我才
        和你算账!他们两个为什么苦呢？}
}
{
    \eA{The sight of their old mistress' delight filled the waiting-maids and
        married women with high glee as well; and each hurried with vehemence to
        execute her respective errand. Those that were to be invited were invited,
        and those that had to be sent for were sent for;}
}
{
}

\Verse{12}
{
    \zA{有了钱也是白填还别人，不如拘了来咱们乐。” 说着早已合了，共凑了一百五十两有零。 贾母道：“一天戏酒用不了。” 尤氏道：
        “既不请客，酒席又不多，两三日的用度都够了。 头等，戏不用钱，省在这上头。” 贾母道：“凤丫头说那一班好，就传那一班。”
        凤姐道：“咱们家的班子都听熟了， 倒是花几个钱叫一班来听听罢。” 贾母道：“这件事我交给珍哥媳妇了，越发叫凤丫 头别操一点心儿，受用一日才算。”
        尤氏答应着。}
}
{
    \eA{and, before the lapse of such time as could suffice to have a meal in, the
        old as well as young, the high as well as low, crammed, in a black mass,
        every bit of the available space in the rooms. Only Mrs. Hsüeh and dowager
        lady Chia sat opposite to each other. Mesdames Hsing and Wang simply seated
        themselves on two chairs, which faced the door of the apartment.}
}
{
}

\Verse{13}
{
    \zA{又说了一回话，都知贾母乏了， 才渐渐的散出来。 尤氏等送出邢夫人王夫人二人散去，因往凤姐房里来，商议怎么办生日的话。
        凤姐儿道：“你不用问我，你只看老太太的眼色儿行事就完了。” 尤氏笑道：“你这 么个阿物儿，也忒行了大运了。 我当有什么事叫我们去，原来单为这个!出了钱不
        算，还叫我操心，你怎么谢我？” 凤姐笑道：“别扯臊!我又没叫你来，谢你什么？ 你怕操心，你这会子就回老太太去，再派一个就是了。”}
}
{
    \eA{Pao-ch'ai and her five or six cousins occupied the stove-couch. Pao-yü sat
        on his grandmother's lap. Below, the whole extent of the floor was crowded
        with inmates on their feet. But old lady Chia forthwith desired that a few
        small stools should be fetched. When brought, these were proffered to Lai
        Ta's mother and some other nurses, who were advanced in years and held in
        respect;}
}
{
}

\Verse{14}
{
    \zA{尤氏笑道：“你瞧瞧，把他 兴的这个样儿!我劝你收着些儿好，太满了就要流出来了。” 二人又说了一回方散。
        次日，将银子送到宁国府来，尤氏方才起来梳洗，因问：“是谁送过来的？” 丫头们回说：“林妈。” 尤氏便命：“叫了他来。” 丫头们走至下房，叫了林之孝家的
        过来。 尤氏命他脚踏上坐了，一面忙着梳洗，一面问他：“这一包银子共多少？” 林之孝家的回说：“这是我们底下人的银子，凑了先送过来。 老太太和太太们的还
        没有呢。”}
}
{
    \eA{for it was the custom in the Chia mansion that the family servants, who had
        waited upon any of the fathers or mothers, should enjoy a higher status than
        even young masters and mistresses. Hence it was that while Mrs. Yu, lady
        Feng and other ladies remained standing below, Lai Ta's mother and three or
        four other old nurses had, after excusing themselves for their rudeness,
        seated themselves on small stools.}
}
{
}

\Verse{15}
{
    \zA{正说着，丫头们回说：“那府里的姨太太打发人送了分子来了。” 尤氏笑 骂道：“小蹄子们，专会记得这些没要紧的话!昨儿不过是老太太一时高兴，故意儿
        的学那小家子凑分子，你们就记得了，到了你们嘴里当正经话说。 还不快接进来呢！” 丫头们笑着忙接银子进来，一共两封，连宝钗、黛玉的都有了。
        尤氏问：“还少谁 的？” 林之孝家的道：“还少老太太、太太、姑娘们的，我们底下姑娘们的。” 尤氏 道：“还有你们大奶奶的呢？”}
}
{
    \eA{Dowager lady Chia recounted, with a face beaming with smiles, the
        suggestions she had shortly made, for the benefit of the various inmates
        present; and one and all, of course, were only too ready to contribute for
        the entertainment. More, some of them, were on friendly terms with lady
        Feng, so they, of their own free will, adopted the proposal;}
}
{
}

\Verse{16}
{
    \zA{林之孝家的道：“奶奶过去，这银子都从二奶奶手里 发，一共都有了。” 说着，尤氏梳洗了，命人伺候车辆。 一时来至荣府，先来见凤姐，只见凤姐已
        将银子封好，正要送去。 尤氏问：“都齐了么？” 凤姐笑道：“都有了!快拿去罢， 丢了我不管。” 尤氏笑道：“我有些信不及，倒要当面点一点。”
        说着，果然按数一 点，只没有李纨的一分。 尤氏笑道：“我说你闹鬼呢!怎么你大嫂子的没有？” 凤姐 笑道：“那么些还不够？ 就短一分儿也罢了。}
}
{
    \eA{others lived in fear and trembling of lady Feng, and these were only too
        anxious to make up to her. Every one, besides, could well afford the means,
        so that, as soon as they heard of the proposed subscriptions, they, with one
        consent, signified their acquiescence. "I'll give twenty taels!" old lady
        Chia was the first to say with a smile playing round her lips.}
}
{
}

\Verse{17}
{
    \zA{等不够了，我再找给你。” 尤氏道：“昨 儿你在人跟前做情，今儿又来和我赖，这我可不依你。 我只和老太太要去。” 凤姐
        笑道：“我看你利害，明儿有了事，我也丁是丁卯是卯的，你也别抱怨！” 尤氏笑道： “只这一分儿不给也罢了，要不看你素日孝敬我，我本来依你么？”
        说着，把平儿 的一分也拿出来，说道：“平儿来把你的收了去，等不够了，我替你添上。” 平儿会 意，笑道：“奶奶先使着，若剩下了，再赏我一样。”}
}
{
    \eA{"I'll follow your lead, dear senior," Mrs. Hsüeh smiled, "and also subscribe
        twenty taels." "We don't presume to place ourselves on an equal footing with
        your ladyship," Mesdames Hsing and Wang pleaded. "We, of course, come one
        degree lower; each of us therefore will contribute sixteen taels." "We too
        naturally rank one step lower," Mrs. Yu and Li Wan also smiled, "so we'll
        each give twelve taels."}
}
{
}

\Verse{18}
{
    \zA{尤氏笑道：“只许你主子作弊， 就不许我作情吗？” 平儿只得收了。 尤氏又道：“我看着你主子这么细致，弄这些 钱，那里使去？
        使不了，明儿带了棺材里使去！” 一面说着，一面又往贾母处来。 先 请了安，大概说了两句话，便走到鸳鸯房中，和鸳鸯商议，只听鸳鸯的主意行事，
        何以讨贾母喜欢。 二人计议妥当。 尤氏临走时，也把鸳鸯的二两银子还他，说：“这 还使不了呢。”}
}
{
    \eA{"You're a widow," dowager lady Chia eagerly demurred, addressing herself to
        Li Wan, "and have lost all your estate, so how could we drag you into all
        this outlay! I'll contribute for you!" "Don't be in such high feather dear
        senior," lady Feng hastily observed laughing, "but just look to your
        accounts before you saddle yourself with this burden! You've already taken
        upon yourself two portions;}
}
{
}

\Verse{19}
{
    \zA{说着，一径出来，又至王夫人跟前说了一回话，因王夫人进了佛堂， 把彩云的一分也还了他。 凤姐儿不在跟前，一时把周赵二人的也还了。 他两个还不
        敢收，尤氏道：“你们可怜见的，那里有这些闲钱？ 凤丫头便知道了，有我应着呢。” 二人听说，千恩万谢的收了。
        转眼已是九月初二日，园中人都打听得尤氏办得十分热闹，不但有戏，连耍百 戏并说书的女先儿全有，都打点着取乐玩耍。 李纨又向众姐妹道：“今儿是正经社
        日，可别忘了。}
}
{
    \eA{and do you now also volunteer sixteen taels on behalf of my elder
        sister-in-law? You may willingly do so, while you speak in the abundance of
        your spirits, but when you, by and bye, come to ponder over what you've
        done, you'll feel sore at heart again! 'It's all that girl Feng that's
        driven me to spend the money,' you'll say in a little time;}
}
{
}

\Verse{20}
{
    \zA{宝玉也不来，想必他不知，又贪住什么玩意儿，把这事又忘了。” 说着，便命丫头：“去瞧做什么呢，快请了来。” 丫头去了半日，回说：“花大姐姐
        说，今儿一早就出门去了。” 众人听了都诧异，说：“再没有出门之理。 这丫头糊涂！” 因又命翠墨去。
        一时翠墨回来，说：“可不真出门了!说有个朋友死了，出去探丧去 了。” 探春道：“断然没有的事。 凭他什么，再没有今日出门之理。 你叫袭人来，我 问他。”}
}
{
    \eA{and you'll devise some ingenious way to inveigle me to fork out three or
        four times as much as your share and thus make up your deficit in an
        underhand way; while I will still be as much in the clouds as if I were in a
        dream!" These words made every one laugh. "According to you, what should be
        done?" dowager lady Chia laughingly inquired. "My birthday hasn't yet come,"
        lady Feng smiled;}
}
{
}

\Verse{21}
{
    \zA{刚说着，只见袭人走来，李纨等都说道：“今儿凭他有什么事，也不该出门。 头一件，你二奶奶的生日，老太太都这么高兴，两府上下都凑热闹儿，他倒走了？
        第二件，又是头一社的正日子，也不告假，就私自去了！” 袭人叹道：“昨儿晚上就 说了，今儿一早有要紧的事，到北静王府里去，就赶着回来。
        劝他别去，他必不依。 今儿一早起来，又要素衣裳穿，想必是北静王府里要紧的什么人没了也未可知。” 李纨等道：“若果如此，也该去走走，只是也该回来了。”}
}
{
    \eA{"and already now I've been the recipient of so much more than I deserve that
        I am quite unhappy. But if I don't contribute a single cash, I shall feel
        really ill at ease for the trouble I shall be giving such a lot of people.
        It would be as well, therefore, that I should bear this share of my senior
        sister-in-law; and, when the day comes, I can eat a few more things, and
        thus be able to enjoy some happiness."}
}
{
}

\Verse{22}
{
    \zA{说着，大家又商议：“咱 们只管作诗，等他来罚他。” 刚说着，只见贾母已打发人来请，便都往前头去了。 袭人回明宝玉的事，贾母不乐，便命人接去。
        原来宝玉心里有件心事，于头一日就吩咐焙茗：“明日一早出门，备两匹马在 后门口等着，不用别人跟着。 说给李贵：我往北府里去了，倘或要有人找我，叫他
        拦住不用找。 只说北府里留下了，横竖就来的。” 焙茗也摸不着头脑，只得依言说 了，今儿一早果然备了两匹马，在园后门等着。}
}
{
    \eA{"Quite right!" cried Madame Hsing and the others at this suggestion. So old
        lady Chia then signified her approval. "There's something more I'd like to
        add," lady Feng pursued smiling.}
}
{
}

\Verse{23}
{
    \zA{天亮了，只见宝玉遍体纯素，从角 门出来，一语不发跨上马，一弯腰顺着街就蹭下去了。 焙茗也只得跨上马，加鞭赶 上，在后面忙问：“往那里去？”
        宝玉道：“这条路是往那里去的？” 焙茗道：“这 是出北门的大道。 出去了冷清清，没有什么玩的。” 宝玉听说，点头道：“正要冷清 清的地方。”
        说着，越发加了两鞭，那马早已转了两个弯子，出了城门。 焙茗越发 不得主意，只得紧紧的跟着。}
}
{
    \eA{"I think that it's fair enough that you, worthy ancestor, should, besides
        your own twenty taels, have to stand two shares as well, the one for cousin
        Liu, the other for cousin Pao-yü, and that Mrs. Hsüeh should, beyond her own
        twenty taels, likewise bear cousin Pao-ch'ai's portion.}
}
{
}

\Verse{24}
{
    \zA{一气跑了七八里路出来，人烟渐渐稀少，宝玉方勒住马，回头问焙茗道：“这 里可有卖香的？” 焙茗道：“香倒有，不知是那一样？” 宝玉想到别的香不好，须
        得檀、芸、降三样。 焙茗笑道：“这三样可难得。” 宝玉为难。 焙茗见他为难，因问 道：“要香做什么使？ 我见二爷时常带的小荷包儿有散香，何不找找？”
        一句提醒了 宝玉，便回手衣襟上挂着个荷包摸了一摸，竟有两星沉速，心内喜欢：“只是不恭 些。” 再想：“自己亲身带的，倒比买的又好些。”}
}
{
    \eA{But it's somewhat unfair that the two ladies Mesdames Hsing and Wang should
        each only give sixteen taels, when their share is small, and when they don't
        subscribe anything for any one else. It's you, venerable senior, who'll be
        the sufferer by this arrangement." Dowager lady Chia, at these words, burst
        out into a boisterous fit of laughter. "It's this hussey Feng," she
        observed, "who, after all, takes my side!}
}
{
}

\Verse{25}
{
    \zA{于是又问炉炭，焙茗道：“这可罢 了，荒郊野外，那里有？ 既用这些，何不早说，带了来岂不便宜？” 宝玉道：“糊涂
        东西!要可以带了来，又不这样没命的跑了。” 焙茗想了半日，笑道：“我得了个主 意，不知二爷心下如何。 我想来二爷不止用这个，只怕还要用别的，这也不是事。
        如今我们索性往前再走二里，就是水仙庵了。” 宝玉听了，忙问：“水仙庵就在这里？ 更好了。 我们就去。”}
}
{
    \eA{What you say is quite right. Hadn't it been for you, I would again have been
        duped by them!" "Dear senior!" lady Feng smiled. Just hand over our two
        cousins to those two ladies and let each take one under her charge and
        finish. If you make each contribute one share, it will be square enough."
        "This is perfectly fair," eagerly rejoined old lady Chia. "Let this
        suggestion be carried out!"}
}
{
}

\Verse{26}
{
    \zA{说着就加鞭前行，一面回头向焙茗道：“这水仙庵的姑子长往 咱们家去，这一去到那里和他借香炉使使，他自然是肯的。” 焙茗道：“别说是咱们
        家的香火，就是平白不认识的庙里，和他借，他也不敢驳回。 只是一件，我常见二 爷最厌这水仙庵的，如何今儿又这样喜欢了？” 宝玉道：“我素日最恨俗人不知原
        故混供神，混盖庙。 这都是当日有钱的老公们和那些有钱的愚妇们，听见有个神， 就盖起庙来供着，也不知那神是何人，因听些野史小说便信真了。}
}
{
    \eA{Lai Ta's mother hastily stood up. "This is such a subversion of right," she
        smiled, "that I'll put my back up on account of the two ladies. She's a
        son's wife, on the other side, and, in here, only a wife's brother's child;
        and yet she doesn't incline towards her mother-in-law and her aunt, but
        takes other people's part. This son's wife has therefore become a perfect
        stranger;}
}
{
}

\Verse{27}
{
    \zA{比如这水仙庵里 面，因供的是洛神，故名水仙庵。 殊不知古来并没有个洛神，那原是曹子建的谎话， 谁知这起愚人就塑了像供着。
        今儿却合我的心事，故借他一用。” 说着，早已来至门前。 那老姑子见宝玉来了，事出意外，竟像天上掉下个活龙 来的一般，忙上来问好，命老道来接马。
        宝玉进去，也不拜洛神之像，却只管赏鉴。 虽是泥塑的，却真有那“翩若惊鸿，婉若游龙”、“荷出渌波，日映朝霞”的姿态。 宝玉不觉滴下泪来。
        老姑子献了茶，宝玉因和他借香炉烧香。}
}
{
    \eA{and a close niece has, in fact, become a distant niece!" As she said this,
        dowager lady Chia and every one present began to laugh. "If the junior
        ladies subscribe twelve taels each," Lai Ta's mother went on to ask, "we
        must, as a matter of course, also come one degree lower; eh?" Upon hearing
        this, old lady Chia remonstrated. "This won't do!" she observed.}
}
{
}

\Verse{28}
{
    \zA{那姑子去了半日，连 香供纸马都预备了来。 宝玉说道：“一概不用。” 命焙茗捧着炉出至后园中，拣一块 干净地方儿，竟拣不出。 焙茗道：“那井台上如何？”
        宝玉点头。 一齐来至井台上，将炉放下，焙茗站过一旁。 宝玉掏出香来焚上，含泪施了半 礼，回身命收了去。
        焙茗答应，且不收，忙爬下磕了几个头，口内祝道：“我焙茗 跟二爷这几年，二爷的心事我没有不知道的，只有今儿这一祭祀，没有告诉我，我 也不敢问。}
}
{
    \eA{"You naturally should rank one degree lower, but you're all, I am well
        aware, wealthy people; and, in spite of your status being somewhat lower,
        your funds are more flourishing than theirs. It's only just then that you
        should be placed on the same standing as those people!" The posse of nurses
        expressed with promptness their acceptance of the proposal their old
        mistress made.}
}
{
}

\Verse{29}
{
    \zA{只是受祭的阴魂，虽不知名姓，想来自然是那人间有一、天上无双，极 聪明清雅的一位姐姐妹妹了。 二爷的心事难出口，我替二爷祝赞你：你若有灵有圣，
        我们二爷这样想着你，你也时常来望候望候二爷，未尝不可。 你在阴间，保佑二爷 来生也变个女孩儿，和你们一处玩耍，岂不两下里都有趣了。” 说毕又磕了几个头，
        才爬起来。 宝玉听他没说完，便掌不住笑了。 因踢他道：“别胡说，看人听见笑话。”}
}
{
    \eA{"The young ladies," dowager lady Chia resumed, "should merely give something
        for the sake of appearances! If each one contributes a sum proportionate to
        her monthly allowance, it will be ample!" Turning her head, "Yüan Yang!" she
        cried, "a few of you should assemble in like manner, and consult as to what
        share you should take in the matter. So bring them along!"}
}
{
}

\Verse{30}
{
    \zA{焙茗 起来，收过香炉，和宝玉走着，因道：“我已经合姑子说了二爷还没用饭，叫他收 拾了些东西，二爷勉强吃些。
        我知道今儿里头大排筵宴，热闹非常，二爷为此才躲 了来的。 横竖在这里清净一天，也就尽乐了； 要不吃东西，断使不得。” 宝玉道：“戏
        酒不吃，这随便的吃些也不妨。” 焙茗道：“这才是。 还有一说：咱们来了，必有人 不放心。 若没有人不放心，便晚些进城何妨？
        若有人不放心，二爷须得进城回家去 才是。}
}
{
    \eA{Yüan Yang assured her that her desires would be duly attended to and walked
        away. But she had not been absent for any length of time, when she appeared
        on the scene along with P'ing Erh, Hsi Jen, Ts'ai Hsia and other girls, and
        a number of waiting-maids as well. Of these, some subscribed two taels;
        others contributed one tael.}
}
{
}

\Verse{31}
{
    \zA{第一老太太、太太也放了心，第二礼也尽了，不过这么着。 就是家去听戏喝 酒，也并不是爷有意，原是陪着父母尽个孝道儿。 要单为这个，不顾老太太、太太
        悬心，就是才受祭的阴魂儿也不安哪。 二爷想我这话怎么样？” 宝玉笑道：“你的 意思我猜着了。 你想着只你一个跟了我出来，回来你怕担不是，所以拿这大题目来
        劝我。 我才来了，不过为尽个礼，再去吃酒看戏，并没说一日不进城。 这已经完了 心愿，赶着进城，大家放心就是了。” 焙茗道：“这更好。”}
}
{
    \eA{"Can it be," dowager lady Chia then said to P'ing Erh, "that you don't want
        any birthday celebrated for your mistress, that you don't range yourself
        also among them?" "The other money I gave," P'ing Erh smiled, "I gave
        privately, and is extra." "This is what I am publicly bound to contribute
        along with the lot." "That's a good child!" lady Chia laughingly rejoined.}
}
{
}

\Verse{32}
{
    \zA{说着二人来至禅堂，果然那姑子收拾了一桌好素菜。 宝玉胡乱吃了些，焙茗也 吃了。 二人便上马，仍回旧路。 焙茗在后面，只嘱咐：“二爷好生骑着。 这马总没
        大骑，手提紧着些儿。” 一面说着，早已进了城，仍从后门进去，忙忙来至怡红院 中。 袭人等都不在屋里，只有几个老婆子看屋子，见他来了，都喜的眉开眼笑道：
        “阿弥陀佛，可来了!没把花姑娘急疯了呢。 上头正坐席呢，二爷快去罢。”}
}
{
    \eA{"Those above as well as those below have all alike given their share," lady
        Feng went on to observe with a smile. "But there are still those two
        secondary wives; are they to give anything or not? Do go and ask them! It's
        but right that we should go to the extreme length and include them.
        Otherwise, they'll imagine that we've looked down upon them!" "Just so!"
        eagerly answered lady Chia, at these words.}
}
{
}

\Verse{33}
{
    \zA{宝玉听 说，忙将素衣脱了，自己找了颜色吉服换上，便问道：“都在什么地方坐席呢？” 老婆子们回道：“在新盖的大花厅上呢。”
        宝玉听了，一径往花厅上来，耳内早隐隐闻得箫管歌吹之声。 刚到穿堂那边， 只见玉钏儿独坐在廊檐下垂泪，一见宝玉来了，便长出了一口气，砸着嘴儿说道：
        “嗳!凤凰来了，快进去罢!再一会子不来，可就都反了。” 宝玉陪笑道：“你猜我往 那里去了？” 玉钏儿把身一扭，也不理他，只管拭泪，宝玉只得怏怏的进去了。}
}
{
    \eA{"How is it that we forgot all about them? The only thing is, I fear, they've
        got no time to spare; yet, tell a servant-girl to go and ask them what
        they'll do!" While she spoke, a servant-girl went off. After a long absence,
        she returned. "Each of them," she reported, "will likewise contribute two
        taels." Dowager lady Chia was delighted with the result.}
}
{
}

\Verse{34}
{
    \zA{到 了花厅上，见了贾母、王夫人等，众人真如得了“凤凰”一般。 贾母先问道：“你 往那里去了，这早晚才来？ 还不给你姐姐行礼去呢！”
        因笑着又向凤姐儿道：“你兄 弟不知好歹，就有要紧的事，怎么也不说一声儿就私自跑了，这还了得!明儿再这 样，等你老子回家，必告诉他打你。”
        凤姐儿笑着道：“行礼倒是小事，宝兄弟明儿 断不可不言语一声儿，也不传人跟着就出去。 街上车马多，头一件叫人不放心。 再 也不像咱们这样人家出门的规矩。”}
}
{
    \eA{"Fetch a pen and inkslab," she cried, "and let's calculate how much they
        amount to, all together." Mrs. Yu abused lady Feng in a low tone of voice.
        "I'll take you, you mean covetous creature, and … ! All these mothers-in-law
        and sisters-in-law have come forward and raised money to celebrate your
        birthday, and are you yet not satisfied that you must also drag in those two
        miserable beings! But what do you do it for?"}
}
{
}

\Verse{35}
{
    \zA{这里贾母又骂跟的人：“为什么都听他的话，说 往那里去就去了，也不回一声儿！” 一面又问：“他到底往那里去了？ 可吃了什么没 有？ 唬着了没有？”
        宝玉只回说：“北静王的一个爱妾没了，今日给他道恼去。 我见 他哭的那样，不好撇下他就回来，所以多等了会子。” 贾母道：“以后再私自出门，
        不先告诉我，一定叫你老子打你！” 宝玉连忙答应着。 贾母又要打跟的人。}
}
{
    \eA{"Try and talk less trash!" lady Feng smiled; also in an undertone. "We'll be
        leaving this place in a little time and then I'll square up accounts with
        you! But why ever are those two miserable? When they have money, they
        uselessly give it to other people; and isn't it better that we should get
        hold of it, and enjoy ourselves with it?"}
}
{
}

\Verse{36}
{
    \zA{众人又 劝道：“老太太也不必生气了，他已经答应不敢了，况且回来又没事，大家该放心 乐一会子了。” 贾母先不放心，自然着急发狠；
        今见宝玉回来，喜且有余，那里还 恨？ 也就不提了。 还怕他不受用，或者别处没吃饭，路上着了惊恐，反又百般的哄 他。 袭人早已过来伏侍，大家仍旧听戏。
        当日演的是《荆钗记》，贾母薛姨妈等都看的心酸落泪，也有笑的，也有恨的， 也有骂的。 要知端底，下回分解。 (本章完)}
}
{
    \eA{While she uttered these taunts, they computed that the collections would
        reach a sum over and above one hundred and fifty taels. "We couldn't
        possibly run through all this for a day's theatricals and banquet!" old lady
        Chia exclaimed. "As no outside guests are to be invited," Mrs. Yu
        interposed, "and the number of tables won't also be many, there will be
        enough to cover two or three days' outlay!}
}
{
}

\Verse{37}
{
    \zA{}
}
{
    \eA{First of all, there won't be anything to spend for theatricals, so we'll
        effect a saving on that item." "Just call whatever troupe that girl Feng may
        say she likes best," dowager lady Chia suggested. "We've heard quite enough
        of the performances of that company of ours," lady Feng said; "let's
        therefore spend a little money and send for another, and see what they can
        do." "I leave that to you, brother Chen's wife," old lady Chia pursued, "in
        order that our girl Feng should have occasion to trouble her mind with as
        little as possible, and be able to enjoy a day's peace and quiet. It's only
        right that she should." Mrs. Yu replied that she would be only too glad to
        do what she could.}
}
{
}

\Verse{38}
{
    \zA{}
}
{
    \eA{They then prolonged their chat for a little longer, until one and all
        realised that their old senior must be quite fagged out, and they gradually
        dispersed. After seeing Mesdames Hsing and Wang off, Mrs. Yu and the other
        ladies adjourned into lady Feng's rooms to consult with her about the
        birthday festivities. "Don't ask me!" lady Feng urged. "Do whatever will
        please our worthy ancestor." "What a fine thing you are to come across such
        a mighty piece of luck!" Mrs. Yu smiled. "I was wondering what had happened
        that she summoned us all! Why, was it simply on this account? Not to breathe
        a word about the money that I'll have to contribute, must I have trouble and
        annoyance to bear as well? How will you show me any thanks?" "Don't bring
        shame upon yourself!" lady Feng laughed. "I didn't send for you;}
}
{
}

\Verse{39}
{
    \zA{}
}
{
    \eA{so why should I be thankful to you! If you funk the exertion, go at once and
        let our venerable senior know, and she'll depute some one else and have
        done." "You go on like this as you see her in such excellent spirits, that's
        why!" Mrs. Yu smilingly answered. "It would be well, I advise you, to pull
        in a bit; for if you be too full of yourself, you'll get your due reward!"
        After some further colloquy, these two ladies eventually parted company. On
        the next day, the money was sent over to the Ning Kuo Mansion at the very
        moment that Mrs. Yu had got up, and was performing her toilette and
        ablutions. "Who brought it?" she asked. "Nurse Lin," the servant-girl said
        by way of response. "Call her in," Mrs. Yu said.}
}
{
}

\Verse{40}
{
    \zA{}
}
{
    \eA{The servant-girls walked as far as the lower rooms and called Lin
        Chih-hsiao's wife to come in. Mrs. Yu bade her seat herself on the
        footstool. While she hurriedly combed her hair and washed her face and
        hands, she wanted to know how much the bundle contained in all. "This is
        what's subscribed by us servants." Lin Chih-hsiao's wife replied, "and so I
        collected it and brought it over first. As for the contributions of our
        venerable mistress, and those of the ladies, they aren't ready yet." But
        simultaneously with this reply, the waiting-maids announced: "Our lady of
        the other mansion and Mrs. Hsüeh have sent over some one with their
        portions." "You mean wenches!" Mrs. Yu cried, scolding them with a smile.
        "All the gumption you've got is to simply bear in mind this sort of
        nonsense!}
}
{
}

\Verse{41}
{
    \zA{}
}
{
    \eA{In a fit of good cheer, your old mistress yesterday purposely expressed a
        wish to imitate those poor people, and raise a subscription. But you at once
        treasured it up in your memory, and, when the thing came to be canvassed by
        you, you treated it in real earnest! Don't you yet quick bundle yourselves
        out, and bring the money in! Be careful and give them some tea before you
        see them off." The waiting-maids smilingly hastened to go and take delivery
        of the money and bring it in. It consisted, in all, of two bundles, and
        contained Pao-ch'ai's and Tai-yü's shares as well. "Whose shares are
        wanting?" Mrs. Yu asked. "Those of our old lady, of Madame Wang, the young
        ladies, and of our girls below are still missing," Lin Chih-hsiao's wife
        explained. "There's also that of your senior lady," Mrs. Yu proceeded.}
}
{
}

\Verse{42}
{
    \zA{}
}
{
    \eA{"You'd better hurry over, my lady," Lin Chih-hsiao's wife said; "for as this
        money will be issued through our mistress Secunda, she'll nobble the whole
        of it." While conversing, Mrs. Yu finished arranging her coiffure and
        performing her ablutions; and, giving orders to see that the carriage was
        got ready, she shortly arrived at the Jung mansion. First and foremost she
        called on lady Feng. Lady Feng, she discovered, had already put the money
        into a packet, and was on the point of sending it over. "Is it all there?"
        Mrs. Yu asked. "Yes, it is," lady Feng smiled, "so you might as well take it
        away at once; for if it gets mislaid, I've nothing to do with it." "I'm
        somewhat distrustful," Mrs. Yu laughed, "so I'd like to check it in your
        presence." These words over, she verily checked sum after sum.}
}
{
}

\Verse{43}
{
    \zA{}
}
{
    \eA{She found Li Wan's share alone wanting. "I said that you were up to tricks!"
        laughingly observed Mrs. Yu. "How is it that your elder sister-in-law's
        isn't here?" "There's all that money; and isn't it yet enough?" lady Feng
        smiled. "If there's merely a portion short it shouldn't matter! Should the
        money prove insufficient, I can then look you up, and give it to you." "When
        the others were present yesterday," Mrs. Yu pursued, "you were ready enough
        to act as any human being would; but here you're again to-day prevaricating
        with me! I won't, by any manner of means, agree to this proposal of yours!
        I'll simply go and ask for the money of our venerable senior." "I see how
        dreadful you are!" lady Feng laughed. "But when something turns up by and
        bye, I'll also be very punctilious;}
}
{
}

\Verse{44}
{
    \zA{}
}
{
    \eA{so don't you then bear me a grudge!" "Well, never mind if you don't give
        your quota!" Mrs. Yu smilingly rejoined. "Were it not that I consider the
        dutiful attentions you've all along shown me would I ever be ready to humour
        you?" So rejoining, she produced P'ing Erh's share. "P'ing Erh, come here,"
        she cried, "take this share of yours and put it away! Should the money
        collected turn out to be below what's absolutely required, I'll make up the
        sum for you." P'ing Erh apprehended her meaning. "My lady," she answered,
        with a cheerful countenance, "it would come to the same thing if you were to
        first spend what you want and to give me afterwards any balance that may
        remain of it."}
}
{
}

\Verse{45}
{
    \zA{}
}
{
    \eA{"Is your mistress alone to be allowed to do dishonest acts," Mrs. Yu
        laughed, "and am I not to be free to bestow a favour?" P'ing Erh had no
        option, but to retain her portion. "I want to see," Mrs. Yu added, "where
        your mistress, who is so extremely careful, will run through all the money,
        we've raised! If she can't spend it, why she'll take it along with her in
        her coffin, and make use of it there." While still speaking, she started on
        her way to dowager lady Chia's suite of rooms. After first paying her
        respects to her, she made a few general remarks, and then betook herself
        into Yüan Yang's quarters where she held a consultation with Yüan Yang.
        Lending a patient ear to all that Yüan Yang;}
}
{
}

\Verse{46}
{
    \zA{}
}
{
    \eA{had to recommend in the way of a programme, and as to how best to give
        pleasure to old lady Chia, she deliberated with her until they arrived at a
        satisfactory decision. When the time came for Mrs. Yu to go, she took the
        two taels, contributed by Yüan Yang, and gave them back to her. "There's no
        use for these!" she said, and with these words still on her lips, she
        straightway quitted her presence and went in search of Madame Wang. After a
        short chat, Madame Wang stepped into the family shrine reserved for the
        worship of Buddha, so she likewise restored Ts'ai Yün's share to her; and,
        availing herself of lady Feng's absence, she presently reimbursed to Mrs.
        Chu and Mrs. Chao the amount of their respective contributions.}
}
{
}

\Verse{47}
{
    \zA{}
}
{
    \eA{These two dames would not however presume to take their money back. "Your
        lot, ladies, is a pitiful one!" Mrs. Yu then expostulated. "How can you
        afford all this spare money! That hussey Feng is well aware of the fact. I'm
        here to answer for you!" At these assurances, both put the money away, with
        profuse expressions of gratitude. In a twinkle, the second day of the ninth
        moon arrived. The inmates of the garden came to find out that Mrs. Yu was
        making preparations on an extremely grand scale; for not only was there to
        be a theatrical performance, but jugglers and women storytellers as well;
        and they combined in getting everything ready that could conduce to afford
        amusement and enjoyment.}
}
{
}

\Verse{48}
{
    \zA{}
}
{
    \eA{"This is," Li Wan went on to say to the young ladies, "the proper day for
        our literary gathering, so don't forget it. If Pao-yü hasn't appeared, it
        must, I presume, be that his mind is so preoccupied with the fuss that's
        going on that he has lost sight of all pure and refined things." Speaking,
        "Go and see what he is up to!" she enjoined a waiting-maid; "and be quick
        and tell him to come." The waiting-maid returned after a long absence.
        "Sister Hua says," she reported, "that he went out of doors, soon after
        daylight this morning." The result of the inquiries filled every one with
        surprise. "He can't have gone out!" they said. "This girl is stupid, and
        doesn't know how to speak." They consequently also directed Ts'ui Mo to go
        and ascertain the truth. In a little time, Ts'ui Mo returned.}
}
{
}

\Verse{49}
{
    \zA{}
}
{
    \eA{"It's really true," she explained, "that he has gone out of doors. He gave
        out that a friend of his was dead, and that he was going to pay a visit of
        condolence." "There's certainly nothing of the kind," T'an Ch'un interposed.
        "But whatever there might have been to call him away, it wasn't right of him
        to go out on an occasion like the present one! Just call Hsi Jen here, and
        let me ask her!" But just as she was issuing these directions, she perceived
        Hsi Jen appear on the scene. "No matter what he may have had to attend to
        to-day," Li Wan and the rest remarked, "he shouldn't have gone out! In the
        first place, it's your mistress Secunda's birthday, and our dowager lady is
        in such buoyant spirits that the various inmates, whether high or low, are
        coming from either mansion to join in the fun;}
}
{
}

\Verse{50}
{
    \zA{}
}
{
    \eA{and lo, he goes off! Secondly, this is the proper day as well for holding
        our first literary gathering, and he doesn't so as apply for leave, but
        stealthily sneaks away." Hsi Jen heaved it sigh. "He said last night," she
        explained, "that he had something very important to do this morning; that he
        was going as far as Prince Pei Ching's mansion, but that he would hurry
        back. I advised him not to go; but, of course, he wouldn't listen to me.
        When he got out of bed, at daybreak this morning, he asked for his plain
        clothes and put them on, so, I suppose, some lady of note belonging to the
        household of Prince Pei Ching must have departed this life; but who can
        tell?" "If such be truly the case," Li Wan and her companions exclaimed,
        "it's quite right that he should have gone over for a while;}
}
{
}

\Verse{51}
{
    \zA{}
}
{
    \eA{but he should have taken care to be back in time !" This remark over, they
        resumed their deliberations. "Let's write our verses," they said, "and we
        can fine him on his return." As these words were being spoken, they espied a
        messenger despatched by dowager lady Chia to ask them over, so they at once
        adjourned to the front part of the compound. Hsi Jen then reported to his
        grandmother what Pao-yü had done. Old lady Chia was upset by the news; so
        much so, that she issued immediate orders to a few servants to go and fetch
        him. Pao-yü had, in fact, been brooding over some affair of the heart. A day
        in advance he therefore gave proper injunctions to Pei Ming. "As I shall be
        going out of doors to-morrow at daybreak," he said, "you'd better get ready
        two horses and wait at the back door! No one else need follow as an escort!}
}
{
}

\Verse{52}
{
    \zA{}
}
{
    \eA{Tell Li Kuei that I've gone to the Pei mansion. In the event of any one
        wishing to start in search of me, bid him place every obstacle in the way,
        as all inquiries can well be dispensed with! Let him simply explain that
        I've been detained in the Pei mansion, but that I shall surely be back
        shortly." Pei Ming could not make out head or tail of what he was driving
        at; but he had no alternative than to deliver his message word for word. At
        the first blush of morning of the day appointed, he actually got ready two
        horses and remained in waiting at the back gate. When daylight set in, he
        perceived Pao-yü make his appearance from the side door; got up, from head
        to foot, in a plain suit of clothes. Without uttering a word, he mounted his
        steed;}
}
{
}

\Verse{53}
{
    \zA{}
}
{
    \eA{and stooping his body forward, he proceeded at a quick step on his way down
        the road. Pei Ming had no help but to follow suit; and, springing on his
        horse, he smacked it with his whip, and overtook his master. "Where are we
        off to?" he eagerly inquired, from behind. "Where does this road lead to?"
        Pao-yü asked. "This is the main road leading out of the northern gate." Pei
        Ming replied. "Once out of it, everything is so dull and dreary that there's
        nothing worth seeing!" Pao-yü caught this answer and nodded his head. "I was
        just thinking that a dull and dreary place would be just the thing!" he
        observed. While speaking, he administered his steed two more whacks. The
        horse quickly turned a couple of corners, and trotted out of the city gate.}
}
{
}

\Verse{54}
{
    \zA{}
}
{
    \eA{Pei Ming was more and more at a loss what to think of the whole affair; yet
        his only course was to keep pace closely in his master's track. With one
        gallop, they covered a distance of over seven or eight lis. But it was only
        when human habitations became gradually few and far between that Pao-yü
        ultimately drew up his horse. Turning his head round: "Is there any place
        here," he asked, "where incense is sold?" "Incense!" Pei Ming shouted, "yes,
        there is; but what kind of incense it is I don't know." "All other incense
        is worth nothing," Pao-yü resumed, after a moment's reflection. "We should
        get sandalwood, conifer and cedar, these three." "These three sorts are very
        difficult to get," Pei Ming smiled. Pao-yü was driven to his wits' ends. But
        Pei Ming noticing his dilemma, "What do you want incense for?"}
}
{
}

\Verse{55}
{
    \zA{}
}
{
    \eA{he felt impelled to ask. "Master Secundus, I've often seen you wear a small
        purse, about your person, full of tiny pieces of incense; and why don't you
        see whether you've got it with you?" This allusion was sufficient to suggest
        the idea to Pao-yü's mind. Forthwith, he drew back his hand and felt the
        purse suspended on the lapel of his coat. It really contained two bits of
        'Ch'en Su.' At this discovery, his heart expanded with delight. The only
        thing that (damped his spirits) was the notion that there was a certain want
        of reverence in his proceedings; but, on second consideration, he concluded
        that what he had about him was, after all, considerably superior to any he
        could purchase, and, with alacrity, he went on to inquire about a censer and
        charcoal. "Don't think of such things!"}
}
{
}

\Verse{56}
{
    \zA{}
}
{
    \eA{Pei Ming urged. "Where could they be procured in a deserted and lonely place
        like this? If you needed them, why didn't you speak somewhat sooner, and we
        could have brought them along with us? Would not this have been more
        convenient?" "You stupid thing!" exclaimed Pao-yü. "Had we been able to
        bring them along, we wouldn't have had to run in this way as if for life!"
        Pei Ming indulged in a protracted reverie, after which, he gave a smile.
        "I've thought of something," he cried, "but I wonder what you'll think about
        it, Master Secundus! You don't, I expect, only require these things; you'll
        need others too, I presume. But this isn't the place for them; so let's move
        on at once another couple of lis, when we'll get to the 'Water Spirit'
        monastery."}
}
{
}

\Verse{57}
{
    \zA{}
}
{
    \eA{"Is the 'Water Spirit' monastery in this neighbourhood?" Pao-yü eagerly
        inquired, upon hearing his proposal. "Yes, that would be better; let's press
        forward." With this reply, he touched his horse with his whip. While
        advancing on their way, he turned round. "The nun in this 'Water Spirit'
        monastery," he shouted to Pei Ming, "frequently comes on a visit to our
        house, so that when we now get there and ask her for the loan of a censer,
        she's certain to let us have it."}
}
{
}

\Verse{58}
{
    \zA{}
}
{
    \eA{"Not to mention that that's a place where our family burns incense," Pei
        Ming answered, "she could not dare to raise any objections, to any appeal
        from us for a loan, were she even in a temple quite unknown to us. There's
        only one thing, I've often been struck with the strong dislike you have for
        this 'Water Spirit' monastery, master, and how is that you're now, so
        delighted with the idea of going to it?" "I've all along had the keenest
        contempt for those low-bred persons," Pao-yü rejoined, "who, without knowing
        why or wherefore, foolishly offer sacrifices to the spirits, and needlessly
        have temples erected.}
}
{
}

\Verse{59}
{
    \zA{}
}
{
    \eA{The reason of it all is, that those rich old gentlemen and unsophisticated
        wealthy women, who lived in past days, were only too ready, the moment they
        heard of the presence of a spirit anywhere, to take in hand the erection of
        temples to offer their sacrifices in, without even having the faintest
        notion whose spirits they were. This was because they readily credited as
        gospel-truth such rustic stories and idle tales as chanced to reach their
        ears. Take this place as an example. Offerings are presented in this 'Water
        Spirit' nunnery to the spirit of the 'Lo' stream; hence the name of 'Water
        Spirit' monastery has been given to it. But people really don't know that in
        past days, there was no such thing as a 'Lo' spirit!}
}
{
}

\Verse{60}
{
    \zA{}
}
{
    \eA{These are, indeed, no better than legendary yarns invented by Ts'ao
        Tzu-chien, and who would have thought it, this sort of stupid people have
        put up images of it, to which they offer oblations. It serves, however, my
        purpose to-day, so I'll borrow of her whatever I need to use." While engaged
        in talking, they reached the entrance. The old nun saw Pao-yü arrive, and
        was thoroughly taken aback. So far was this visit beyond her expectations,
        that well did it seem to her as if a live dragon had dropped from the
        heavens. With alacrity, she rushed up to him; and making inquiries after his
        health, she gave orders to an old Taoist to come and take his horse. Pao-yü
        stepped into the temple. But without paying the least homage to the image of
        the 'Lo' spirit, he simply kept his eyes fixed intently on it;}
}
{
}

\Verse{61}
{
    \zA{}
}
{
    \eA{for albeit made of clay, it actually seemed, nevertheless, to flutter as
        does a terror-stricken swan, and to wriggle as a dragon in motion. It looked
        like a lotus, peeping its head out of the green stream, or like the sun,
        pouring its rays upon the russet clouds in the early morn. Pao-yü's tears
        unwittingly trickled down his cheeks. The old nun presented tea. Pao-yü then
        asked her for the loan of a censer to burn incense in. After a protracted
        absence, the old nun returned with some incense as well as several paper
        horses, which she had got ready for him to offer. But Pao-yü would not use
        any of the things she brought. "Take the censer," he said to Pei Ming, "and
        go out into the back garden and find a clean spot!" But having been unable
        to discover one;}
}
{
}

\Verse{62}
{
    \zA{}
}
{
    \eA{"What about, the platform round that well?" Pei Ming inquired. Pao-yü nodded
        his head assentingly. Then along with him, he repaired to the platform of
        the well. He deposited the censer on the ground, while Pei Ming stood on one
        side. Pao-yü produced the incense, and threw it on the fire. With suppressed
        tears, he performed half of the ceremony, and, turning himself round, he
        bade Pei Ming clear the things away. Pei Ming acquiesced; but, instead of
        removing the things, he speedily fell on his face, and made several
        prostrations, as his lips uttered this prayer: "I, Pei Ming, have been in
        the service of Master Secundus for several years. Of the secrets of Mr.
        Secundus' heart there are none, which I have not known, save that with
        regard to this sacrifice to-day; the object of which, he has neither told
        me;}
}
{
}

\Verse{63}
{
    \zA{}
}
{
    \eA{nor have I had the presumption to ask. But thou, oh spirit! who art the
        recipient of these sacrificial offerings, must, I expect, unknown though thy
        surname and name be to me, be a most intelligent and supremely beautiful
        elder or younger sister, unique among mankind, without a peer even in
        heaven! As my Master Secundus cannot give vent to the sentiments, which fill
        his heart, allow me to pray on his behalf! Should thou possess spirituality,
        and holiness be thy share, do thou often come and look up our Mr. Secundus,
        for persistently do his thoughts dwell with thee! And there is no reason why
        thou should'st not come! But should'st thou be in the abode of the dead,
        grant that our Mr. Secundus too may, in his coming existence, be transformed
        into a girl, so that he may be able to amuse himself with you all!}
}
{
}

\Verse{64}
{
    \zA{}
}
{
    \eA{And will not this prove a source of pleasure to both sides?" At the close of
        his invocation, he again knocked his head several times on the ground, and,
        eventually, rose to his feet. Pao-yü lent an ear to his utterances, but,
        before they had been brought to an end, he felt it difficult to repress
        himself from laughing. Giving him a kick, "Don't talk such stuff and
        nonsense!" he shouted. "Were any looker-on to overhear what you say, he'd
        jeer at you!" Pei Ming got up and put the censer away. While he walked along
        with Pao-yü, "I've already," he said, "told the nun that you hadn't as yet
        had anything to eat, Master Secundus, and I bade her get a few things ready
        for you, so you must force yourself to take something.}
}
{
}

\Verse{65}
{
    \zA{}
}
{
    \eA{I know very well that a grand banquet will be spread in our mansion to-day,
        that exceptional bustle will prevail, and that you have, on account of this,
        Sir, come here to get out of the way. But as you're, after all, going to
        spend a whole day in peace and quiet in here, you should try and divert
        yourself as best you can. It won't, therefore, by any manner of means do for
        you to have nothing to eat." "I won't be at the theatrical performance to
        have any wine," Pao-yü remarked, "so what harm will there be in my having a
        drink here, as the fancy takes me?" "Quite so!" rejoined Pei Ming. "But
        there's another consideration. You and I have run over here; but there must
        be some whose minds are ill at ease.}
}
{
}

\Verse{66}
{
    \zA{}
}
{
    \eA{Were there no one uneasy about us, well, what would it matter if we got back
        into town as late as we possibly could? But if there be any solicitous on
        your account, it's but right, Master Secundus, that you should enter the
        city and return home. In the first place, our worthy old mistress and Madame
        Wang, will thus compose their minds; and secondly, you'll observe the proper
        formalities, if you succeed in doing nothing else. But even supposing that,
        when once you get home, you feel no inclination to look at the plays and
        have anything to drink, you can merely wait upon your father and mother, and
        acquit yourself of your filial piety!}
}
{
}

\Verse{67}
{
    \zA{}
}
{
    \eA{Well, if it's only a matter of fulfilling this obligation, and you don't
        care whether our old mistress and our lady, your mother, experience concern
        or not, why, the spirit itself, which has just been the recipient of your
        oblations, won't feel in a happy frame of mind! You'd better therefore,
        master, ponder and see what you think of my words!" "I see what you're
        driving at!" Pao-yü smiled. "You keep before your mind the thought that
        you're the only servant, who has followed me as an attendant out of town,
        and you give way to fear that you will, on your return, have to bear the
        consequences. You hence have recourse to these grandiloquent arguments to
        shove words of counsel down my throat!}
}
{
}

\Verse{68}
{
    \zA{}
}
{
    \eA{I've come here now with the sole object of satisfying certain rites, and
        then going to partake of the banquet and be a spectator of the plays; and I
        never mentioned one single word about any intention on my part not to go
        back to town for a whole day! I've, however, already accomplished the wish I
        fostered in my heart, so if we hurry back to town, so as to enable every one
        to set their solicitude at rest, won't the right principle be carried out to
        the full in one respect as well as another?" "Yes, that would be better!"
        exclaimed Pei Ming. Conversing the while, they wended their way into the
        Buddhistic hall.}
}
{
}

\Verse{69}
{
    \zA{}
}
{
    \eA{Here the nun had, in point of fact, got ready a table with lenten viands.
        Pao-yü hurriedly swallowed some refreshment and so did Pei Ming; after
        which, they mounted their steeds and retraced their steps homewards, by the
        road they had come. Pei Ming followed behind. "Master Secundus!" he kept on
        shouting, "be careful how you ride! That horse hasn't been ridden very much,
        so hold him in tight a bit." As he urged him to be careful, they reached the
        interior of the city walls, and, making their entrance once more into the
        mansion by the back gate, they betook themselves, with all possible
        despatch, into the I Hung court. Hsi Jen and the other maids were not at
        home. Only a few old women were there to look after the rooms.}
}
{
}

\Verse{70}
{
    \zA{}
}
{
    \eA{As soon as they saw him arrive, they were so filled with gratification that
        their eyebrows dilated and their eyes smiled. "O-mi-to-fu!" they said
        laughingly, "you've come! You've all but driven Miss Hua mad from despair!
        In the upper quarters, they're just seated at the feast, so be quick, Mr.
        Secundus, and go and join them." At these words, Pao-yü speedily divested
        himself of his plain clothes and put on a coloured costume, reserved for
        festive occasions, which he hunted up with his own hands. This done, "Where
        are they holding the banquet?" he inquired. "They're in the newly erected
        large reception pavilion," the old women responded. Upon catching their
        reply, Pao-yü straightway started for the reception-pavilion.}
}
{
}

\Verse{71}
{
    \zA{}
}
{
    \eA{From an early moment, the strains of flageolets and pipes, of song and of
        wind-instruments faintly fell on his ear. The moment he reached the passage
        on the opposite side, he discerned Yü Ch'uan-erh seated all alone under the
        eaves of the verandah giving way to tears. As soon as she became conscious
        of Pao-yü's arrival, she drew a long, long breath. Smacking her lips, "Ai!"
        she cried, "the phoenix has alighted! go in at once! Hadn't you come for
        another minute, every one would have been quite upset!" Pao-yü forced a
        smile. "Just try and guess where I've been?" he observed. Yü Ch'uan-erh
        twisted herself round, and, paying no notice to him, she continued drying
        her tears. Pao-yü had, therefore, no option but to enter with hasty step.}
}
{
}

\Verse{72}
{
    \zA{}
}
{
    \eA{On his arrival in the reception-hall, he paid his greetings to his
        grandmother Chia, to Madame Wang, and the other inmates, and one and all
        felt, in fact, as happy to see him back as if they had come into the
        possession of a phoenix. "Where have you been," dowager lady Chia was the
        first to ask, "that you come back at this hour? Don't you yet go and pay
        your congratulations to your cousin?" And smiling she proceeded, addressing
        herself to lady Feng, "Your cousin has no idea of what's right and what's
        wrong. Even though he may have had something pressing to do, why didn't he
        utter just one word, but stealthily bolted away on his own hook? Will this
        sort of thing ever do?}
}
{
}

\Verse{73}
{
    \zA{}
}
{
    \eA{But should you behave again in this fashion by and bye, I shall, when your
        father comes home, feel compelled to tell him to chastise you." Lady Feng
        smiled. "Congratulations are a small matter?" she observed. "But, cousin
        Pao, you must, on no account, sneak away any more without breathing a word
        to any one, and not sending for some people to escort you, for carriages and
        horses throng the streets. First and foremost, you're the means of making
        people uneasy at heart; and, what's more, that isn't the way in which
        members of a family such as ours should go out of doors!" Dowager lady Chia
        meanwhile went on reprimanding the servants, who waited on him. "Why," she
        said, "do you all listen to him and readily go wherever he pleases without
        even reporting a single word? But where did you really go?"}
}
{
}

\Verse{74}
{
    \zA{}
}
{
    \eA{Continuing, she asked, "Did you have anything to eat? Or did you get any
        sort of fright, eh?" "A beloved wife of the duke of Pei Ching departed this
        life," Pao-yü merely returned for answer, "and I went to-day to express my
        condolences to him. I found him in such bitter anguish that I couldn't very
        well leave him and come back immediately. That's the reason why I tarried
        with him a little longer." "If hereafter you do again go out of doors slyly
        and on your own hook," dowager lady Chia impressed on his mind, "without
        first telling me, I shall certainly bid your father give you a caning!"
        Pao-yü signified his obedience with all promptitude. His grandmother Chia
        was then bent upon having the servants, who were on attendance on him,
        beaten, but the various inmates did their best to dissuade her.}
}
{
}

\Verse{75}
{
    \zA{}
}
{
    \eA{"Venerable senior!" they said, "you can well dispense with flying into a
        rage! He has already promised that he won't venture to go out again.
        Besides, he has come back without any misadventure, so we should all compose
        our minds and enjoy ourselves a bit!" Old lady Chia had, at first, been full
        of solicitude. She had, as a matter of course, been in a state of despair
        and displeasure; but, seeing Pao-yü return in safety, she felt immoderately
        delighted, to such a degree, that she could not reconcile herself to visit
        her resentment upon him. She therefore dropped all mention of his escapade
        at once.}
}
{
}

\Verse{76}
{
    \zA{}
}
{
    \eA{And as she entertained fears lest he may have been unhappy or have had, when
        he was away, nothing to eat, or got a start on the road, she did not punish
        him, but had, contrariwise, recourse to every sort of inducement to coax him
        to feel at ease. But Hsi Jen soon came over and attended to his wants, so
        the company once more turned their attention to the theatricals. The play
        acted on that occasion was, "The record of the boxwood hair-pin." Dowager
        lady Chia, Mrs. Hsüeh and the others were deeply impressed by what they saw
        and gave way to tears. Some, however, of the inmates were amused; others
        were provoked to anger; others gave vent to abuse. But, reader, do you wish
        to know the sequel? If so, the next chapter will explain it.}
}
{
}
